\section{The energies of the actual integral maps}
\label{rh-energy}

An arithmetic map has both an integral meaning and a complex matrix
representation. The following estimates simplify the required matrix
conditions. Their use in the arithmetic argument requires the matrices
to represent the same integral maps, in one common target basis, with
the prescribed positive normalizers and metrics.

\subsection{Two aggregate energies}

Let $T_i:H_i\to H$ be finitely many linear maps of finite-dimensional
complex Hermitian spaces, with $r_i=\dim H_i>0$, $D=\dim H>0$ and
$R=\sum_i r_i$. Define
\[
 Q=\sum_iT_iT_i^*,\quad M=\operatorname{tr}Q>0,
 \quad \mathcal S=\bigoplus_iT_i^*T_i,
\]
\begin{equation}\label{rh-energies}
 \Delta=\sum_i\|T_i^*T_i-I_{r_i}\|_{\rm HS}^2,
 \qquad V=\left\|Q-\frac MD I_D\right\|_{\rm HS}^2.
\end{equation}
The Hilbert--Schmidt norm satisfies
$\|A\|_{\rm HS}^2=\operatorname{tr}(A^*A)$.

\begin{lemma}[Exact energy identities]\label{rh-energy-identities}
The diagonal defect and target isotropy defect are
\[
 e=\Delta/R,\qquad
 \eta=\frac{DV}{M^2}=\frac{D\operatorname{tr}(Q^2)}{M^2}-1.
\]
Moreover,
\[
 |M-R|\le\sqrt{R\Delta},\qquad
 \rank Q\ge\frac D{1+\eta},\qquad
 \sum_i\dim\ker T_i\le\Delta.
\]
If $\Omega=\sum_{i\ne j}\|T_i^*T_j\|_{\rm HS}^2$, then
$\operatorname{tr}(Q^2)=\Omega+2M-R+\Delta$.
\end{lemma}

\begin{proof}
Expand the squared norms. Cauchy--Schwarz applied to
$\mathcal S-I_R$ gives the trace inequality. Applied to the
nonzero eigenvalues of $Q$, it gives
$M^2\le\rank(Q)\operatorname{tr}(Q^2)$.
Each vector in the kernel of $T_i$ contributes an eigenvalue $-1$
to $T_i^*T_i-I$. Finally, expand $\operatorname{tr}(Q^2)$ into its
diagonal and off-diagonal terms.
\end{proof}

The diagonal energy alone is insufficient. For $D$ one-dimensional
sources all mapped isometrically to the same line in $\mathbb C^D$,
one has $\Delta=0$ but $\eta=D-1$. A global isotropy estimate is
therefore a genuine additional input.

\begin{proposition}[Compression measured by retained energy]
\label{rh-energy-compression}
Let $J_i:H'_i\to H_i$ be isometries. Put $T'_i=T_iJ_i$ and attach
primes to the corresponding quantities. Suppose $R'>0$ and $M'>0$.
With $l=R-R'$, $L=M-M'$ and $\alpha=M/M'$, one has
\[
 0\le Q'\le Q,\quad \Delta'\le\Delta,\quad
 |L-l|\le\sqrt{l\Delta},\quad
 M'\ge R'-\sqrt{R'\Delta},
\]
\begin{equation}\label{rh-compressed-isotropy}
 \eta'\le
 \left(\frac{\alpha\sqrt\eta+
                \sqrt{\alpha^2\eta+4(\alpha-1)}}2\right)^2.
\end{equation}
For a perfectly isotropic full family, $\eta=0$, this becomes
$\eta'\le L/(M-L)$, and equality is possible.
\end{proposition}

\begin{proof}
Each $J_iJ_i^*$ is an orthogonal projection. This gives $Q'\le Q$;
compression of $T_i^*T_i-I$ contracts its Hilbert--Schmidt norm,
giving $\Delta'\le\Delta$. Taking traces on the discarded and
retained source subspaces and applying Cauchy--Schwarz gives the
two trace bounds.
Since $\operatorname{tr}((Q-Q')Q')\ge0$, centering both matrices gives
\[
 \frac{(M')^2}{D}(1+\eta')
 \le\operatorname{tr}(QQ')
 \le\frac{MM'}D(1+\sqrt\eta\sqrt{\eta'}).
\]
Solve the resulting quadratic inequality in $\sqrt{\eta'}$.
For sharpness when $\eta=0$, take $Q=cI_D$ and $Q'=cP$ for an
orthogonal projection $P$ of positive rank $k$. Then
$\eta'=(D-k)/k=L/(M-L)$.
\end{proof}

The denominators $M'$ and $R'$ are substantive. A bound on the
number of discarded dimensions does not by itself measure the
energy they carry. The exact retained-energy formula may be used
before any uniform estimate for those denominators.

\subsection{A centered comparison of dimensions and slopes}

Define the nonnegative scale-invariant quantities $a,b\ge0$ by
\[
 a^2=\frac R{M^2}\left\|\mathcal S-\frac MR I_R\right\|_{\rm HS}^2
     =\frac{R\Delta-(M-R)^2}{M^2},\qquad b^2=\eta.
\]

\begin{lemma}[Centered dimension comparison]\label{rh-centered-rank}
If subspaces $V_i\subset H_i$ have total dimension $k$ and
$T_i(V_i)\subset W\subset H$, then, with $p=k/R$ and $q=\dim W/D$,
\begin{equation}\label{rh-centered-rank-bound}
               p-a\sqrt{p(1-p)}\le q+b\sqrt{q(1-q)}.
\end{equation}
\end{lemma}

\begin{proof}
Let $P=\bigoplus_iP_{V_i}$ and $C=\sum_iT_iP_{V_i}T_i^*$.
Since $\|P-pI_R\|_{\rm HS}^2=Rp(1-p)$,
\[
 \operatorname{tr}C=\operatorname{tr}(P\mathcal S)
                         \ge M\{p-a\sqrt{p(1-p)}\}.
\]
Also $0\le C\le Q$, and its range lies in $W$. Centering $Q$ gives
\[
 \operatorname{tr}C\le\operatorname{tr}(P_WQ)
                         \le M\{q+b\sqrt{q(1-q)}\}.
\]
Divide by $M>0$.
\end{proof}

For the arithmetic use, fix a number field $K$ and let all lattices
be defined over $K$. Let $\mu_{\max}$ and $\mu_{\min}$ denote
the extreme slopes of the Harder--Narasimhan filtration. Its slopes,
counted with their multiplicities, average to $\mu(E)$. For a nonzero
integral map $f:E\to F$, use the ordinary height
\[
 h(f)=\frac1{[K:\mathbb Q]}\sum_{\sigma:K\hookrightarrow\mathbb C}
                                      \log\|f_\sigma\|;
\]
put $h(0)=-\infty$. A height-nonpositive integral map carries the
source subspace of slopes at least $\xi$ into the corresponding
target subspace. To see this, any nonzero induced map to the target
quotient would have an image of slope at least $\xi$ from the source
and of slope strictly less than $\xi$ from the quotient, contrary
to the determinant inequality. This is the filtration functoriality
used below.

\begin{theorem}[A centered arithmetic slope bound]
\label{rh-centered-slope}
Let $f_i:E_i\to F$ be integral maps, with source ranks $r_i>0$,
target rank $n>0$, and heights $h(f_i)\le h_i$. Suppose
\[
 \frac1R\sum_ir_i\{\mu(E_i)-h_i\}\ge s,\quad
 \mu_{\max}(E_i)-h_i\le s+U,\quad
 \mu_{\min}(F)\ge s-H,\qquad U,H\ge0.
\]
Suppose the first $n\le D$ rows of the matrices $T_i$ represent
these maps at one complex embedding, in individual source bases and
one common target basis. Put $\nu=n/D$, $\delta_0=1-\nu$, and
\[
 \Phi(U,H)=\begin{cases}\sqrt{UH},&U\le H,\\(U+H)/2,&U>H.
 \end{cases}
\]
Then
\begin{equation}\label{rh-centered-slope-bound}
 s-\mu(F)\le P_{\rm cen}:=
       \frac{\delta_0U+a\Phi(U,H)+\sqrt\eta(U+H)/2}{\nu}.
\end{equation}
\end{theorem}

\begin{proof}
Multiply the source norms by $e^{h_i}$. Their slopes decrease by
$h_i$ and all map heights become nonpositive. For
$\xi\in[s-H,s+U]$, let $p(\xi)$ be the total source filtration
rank divided by $R$, and $q(\xi)$ the target filtration rank divided
by $n$. Functoriality puts the images in a subspace of fraction
$w(\xi)=\delta_0+\nu q(\xi)$ of the full target: include its
$D-n$ discarded coordinates. Lemma~\ref{rh-centered-rank} gives
\[
 \nu q\ge p-\delta_0-a\sqrt{p(1-p)}
                            -\sqrt\eta\sqrt{w(1-w)}.
\]
On this interval, the upper source slope bound and its mean give
$\int p\ge H$ and $\int(1-p)\le U$.
If $z=\int(1-p)$ and $L_0=U+H$, Cauchy--Schwarz yields
\[
 \int\sqrt{p(1-p)}\le\sqrt{z(L_0-z)}\le\Phi(U,H),
 \qquad \int\sqrt{w(1-w)}\le L_0/2.
\]
The lower target slope bound gives
$\mu(F)\ge s-H+\int q$. Integrating the preceding dimension
inequality cancels the terms in $\delta_0H$ and proves
\eqref{rh-centered-slope-bound}. If $U+H=0$, the conclusion follows
directly from the lower target slope bound.
\end{proof}

For $a=\eta=\delta_0=0$, this cost is zero. The target-cut term is
sharp: project a full identity matrix on $D$ coordinates to $n$
coordinates, give the discarded source directions slope $s+U$ and
the retained ones slope $s-\delta_0U/\nu$, and use identical metrics
on the latter and the target. The source mean is $s$, the map has
height zero, and the target slope deficit is $\delta_0U/\nu$.
This satisfies the target minimum-slope hypothesis whenever
$H\ge\delta_0U/\nu$.

The centered bound need not be the smallest when $\eta$ is small
and $H$ is large. The GRH manuscript's proved \emph{Refined slope
comparison} \cite{rh-companion}, under the same lattice conditions
and $e<1$, gives
\[
 P_{\rm ref}=\frac{A+B/2+\sqrt{AB+B^2/4}}{\nu},
 \qquad B=\eta H,
\]
\[
 A=\frac{U\log(1+H/U)+2\sqrt e\{\sqrt{U(U+H)}-U\}}{1-\sqrt e}.
\]
At $U=0$ set $A=0$. One may use
$P=\min\{H,P_{\rm cen},P_{\rm ref}\}$ whenever all hypotheses hold.
The first bound follows directly from $\mu(F)\ge s-H$.

An optimal common positive rescaling of the matrices gives
\[
 e_* =\min_{t>0}\frac1R\|t\mathcal S-I_R\|_{\rm HS}^2
       =1-\frac{M^2}{R\operatorname{tr}(\mathcal S^2)}
       =\frac{a^2}{1+a^2}<1.
\]
This follows by completing the square; the minimum is at
$t=M/\operatorname{tr}(\mathcal S^2)$. Absorbing the rescaling into
the source bases preserves the represented arithmetic maps and their
heights, and the target isotropy defect is unchanged. The denominator
$1-\sqrt{e_*}$ must still be retained.

\subsection{A global kernel formulation}

Realize $H$ isometrically in $L^2(X,\mu)$, where $(X,\mu)$ is
sigma-finite, and choose orthonormal bases of $H_i$ and $H$. Let
$\theta_{ij}$ be the images of the source basis vectors. Put
\[
 K_i(z,w)=\sum_j\theta_{ij}(z)\overline{\theta_{ij}(w)},\quad
 K=\sum_iK_i,\quad P_H(z,w)=\sum_{a=1}^Dh_a(z)\overline{h_a(w)}.
\]
Finite expansion and Parseval give
\begin{align*}
 M&=\int K(z,z)\,d\mu(z),\\
 \Delta&=\sum_i\left\{\iint|K_i(z,w)|^2\,d\mu(z)d\mu(w)
                    -2\int K_i(z,z)\,d\mu(z)+r_i\right\},\\
 \eta&=\frac D{M^2}\iint\left|K(z,w)-\frac MD P_H(z,w)\right|^2
                                                   \,d\mu(z)d\mu(w).
\end{align*}
Indeed, $K_i$ represents $T_iT_i^*$ and the $L^2$ norm of its
kernel is its Hilbert--Schmidt norm. The displayed finite sums specify
the representatives used on the diagonal. The centered integral
retains the interference between all maps; a bound on its diagonal
alone is not the required estimate.

For theta maps the necessary input is an identity between these
normalized kernels and the matrices of the actual integral maps.
Every individual normalizer must be positive. An appropriate period
or Bessel identity can supply that identification, but the numerical
global energy estimate and the arithmetic height bounds remain
distinct obligations.

\subsection{A one-sided period bound can supply the target energy}

An alternative to computing every interaction is to prove one upper
Bessel bound for the same common target. It must be calibrated to
the actual total mass. The following two propositions give the
resulting energy estimate and its behavior under both source and
target restriction.

\begin{proposition}[A one-sided period estimate]\label{shared-period-bessel}
Let \(H,H_i\) be finite-dimensional complex Hilbert spaces, indexed by a
finite set, with \(D=\dim H>0\). Let \(P_i:H_i\to H\), and fix real
constants \(c_i>0\). Put
\[
 T_i=c_i^{-1/2}P_i,\qquad Q=\sum_iT_iT_i^*,\qquad
 M=\operatorname{tr}Q>0.
\]
Suppose that, for some real \(\alpha\),
\begin{equation}\label{shared-period-bessel-bound}
 \sum_i c_i^{-1}\|P_i^*h\|^2\le\alpha\|h\|^2
 \qquad(h\in H).
\end{equation}
Then \(\alpha>0\), and the centered target energy satisfies
\begin{equation}\label{shared-period-bessel-energy}
 0\le
 \eta:=\frac D{M^2}
       \left\|Q-\frac MD I_H\right\|_{\mathrm{HS}}^2
 \le\lambda:=\frac{\alpha D}{M}-1.
\end{equation}
Equivalently, the hypothesis is \(Q\le\alpha I_H\), or
\[
 \left\|\sum_iT_i v_i\right\|^2
       \le\alpha\sum_i\|v_i\|^2
       \qquad(v_i\in H_i).
\]
\end{proposition}

\begin{proof}
The first equivalence follows by evaluating the quadratic form of \(Q\).
The second uses equality of the norms of the synthesis map
\(\bigoplus_iH_i\to H\) and its adjoint. Every eigenvalue \(q_a\) of
\(Q\) lies in \([0,\alpha]\), so
\[
 \operatorname{tr}(Q^2)\le\alpha M,\qquad
 \left\|Q-\frac MD I_H\right\|_{\mathrm{HS}}^2
       =\operatorname{tr}(Q^2)-\frac{M^2}{D}.
\]
This proves the energy bound. Taking traces also gives
\(\alpha D\ge M>0\).
\end{proof}

\begin{proposition}[Restriction to the final source and target]
\label{shared-period-joint-restriction}
In Proposition~\ref{shared-period-bessel}, let \(J_i:E_i\to H_i\) be
isometries, let \(F\subset H\) have dimension \(n>0\), and let \(\Pi\)
be its orthogonal projection. Define
\[
 T_i'=T_iJ_i,\quad A_i=\Pi T_i':E_i\to F,\quad
 Q'=\sum_iT_i'T_i'^*,\quad Q_0=\sum_iA_iA_i^*,
\]
\[
 L=\operatorname{tr}(Q-Q'),\quad \rho=L/M,\quad
 \delta=(D-n)/D,\quad
 \kappa=1-\rho-(1+\lambda)\delta.
\]
If \(\kappa>0\), then
\begin{align}
 M_0:=\operatorname{tr}Q_0&\ge M\kappa>0,\label{shared-period-retained-mass}\\
 \eta_0:=
 \frac n{M_0^2}\left\|Q_0-\frac{M_0}{n}I_F\right\|_{\mathrm{HS}}^2
 &\le
 \frac{\lambda+\rho}{1-\rho-(1+\lambda)\delta}.
 \label{shared-period-retained-energy}
\end{align}
Moreover, if
\[
 \Delta=\sum_i\|T_i^*T_i-I_{H_i}\|_{\mathrm{HS}}^2,\qquad
 \Delta_0=\sum_i\|A_i^*A_i-I_{E_i}\|_{\mathrm{HS}}^2,
\]
then
\begin{equation}\label{shared-period-retained-diagonal}
 \Delta_0\le\Delta+2\alpha(D-n).
\end{equation}
If \(l=\sum_i(\dim H_i-\dim E_i)\), one may use
\begin{equation}\label{shared-period-discarded-mass}
 0\le L\le l+\sqrt{l\Delta}.
\end{equation}
\end{proposition}

\begin{proof}
Orthogonal domain restriction gives \(0\le Q'\le Q\le\alpha I_H\).
Consequently
\[
 \operatorname{tr}((I-\Pi)Q')\le\alpha(D-n),\qquad
 M_0=M-L-\operatorname{tr}((I-\Pi)Q')\ge M\kappa.
\]
On \(F\) we have \(Q_0\le\alpha I_F\). The preceding proposition gives
\[
 \eta_0\le\frac{\alpha n}{M_0}-1
 \le\frac{(1+\lambda)(1-\delta)}{\kappa}-1
 =\frac{\lambda+\rho}{\kappa}.
\]
For the diagonal estimate, work on \(\bigoplus_iE_i\) and write
\[
 G=\operatorname{diag}(T_i'^*T_i'),\qquad
 B=\operatorname{diag}(T_i'^*(I-\Pi)T_i').
\]
Then \(0\le B\le G\), the final diagonal Gram operator is \(G-B\),
and domain compression gives
\(\Delta':=\|G-I\|_{\mathrm{HS}}^2\le\Delta\).
Since \(\operatorname{tr}((G-B)B)\ge0\), even without commutation,
\[
 \begin{split}
 \Delta_0
 &=\Delta'-2\operatorname{tr}(GB)+2\operatorname{tr}B+
                          \operatorname{tr}(B^2)\\
 &\le\Delta'+2\operatorname{tr}B-\operatorname{tr}(B^2)
 \le\Delta+2\alpha(D-n).
 \end{split}
\]
Here \(\operatorname{tr}B=\operatorname{tr}((I-\Pi)Q')\).
Finally \(L\) is the trace of
\(\operatorname{diag}(T_i^*T_i)\) on the discarded domain subspace.
Subtract its identity operator and apply Cauchy--Schwarz on that
\(l\)-dimensional subspace to obtain
\eqref{shared-period-discarded-mass}.
\end{proof}

These estimates apply to the actual final maps \(A_i:E_i\to F\).
Their diagonal normalization uses
\(R'=\sum_i\dim E_i\); a frame comparison requiring
\(\sqrt{\Delta_0/R'}<1\) still requires that inequality.
One may use the smaller of \eqref{shared-period-retained-diagonal}
and any other valid diagonal bound. The lattice and height hypotheses
must refer to these same maps.
If \(\lambda=\rho=0\), the target energy stays zero under every
orthogonal target restriction. The upper Bessel estimate must therefore
be calibrated to the actual mass \(M\); a coarse operator bound alone
does not supply a small frame error.

\begin{proposition}[A signed orbital criterion]\label{shared-period-block-gram}
Use the spaces, maps and positive constants of
Proposition~\ref{shared-period-bessel}. Suppose the normalized Gram
operators have absolutely convergent operator-norm expansions
\[
 K_{ij}=c_i^{-1/2}c_j^{-1/2}P_i^*P_j
       =\sum_{t\in\mathcal T_{ij}}a_{ij,t},
 \qquad \sum_t\|a_{ij,t}\|_{\mathrm{op}}<\infty.
\]
For each \(i\), suppose a specified subset
\(\mathcal I_i\subset\mathcal T_{ii}\) has sum \(I_{H_i}\).
On \(\mathcal E=\bigoplus_iH_i\), define \(N=[N_{ij}]\) by
\[
 N_{ii}=\sum_{t\notin\mathcal I_i}a_{ii,t},
 \qquad N_{ij}=\sum_{t\in\mathcal T_{ij}}a_{ij,t}\quad(i\ne j).
\]
If, for every \(v=(v_i)\in\mathcal E\),
\begin{equation}\label{shared-period-signed-orbital-bound}
 \Re\sum_{i,j}\langle v_i,N_{ij}v_j\rangle
       \le(\alpha-1)\sum_i\|v_i\|^2,
\end{equation}
then the upper Bessel estimate
\eqref{shared-period-bessel-bound} holds.
It suffices to establish the same quadratic-form bound, with the same
\(\alpha\), for finite orbital truncations converging in operator norm
to \(N\).
\end{proposition}

\begin{proof}
Let \(S:\mathcal E\to H\) be \(S(v)=\sum_iT_iv_i\).
Absolute convergence and the diagonal identity give
\[
 S^*S=[K_{ij}]=I_{\mathcal E}+N,\qquad SS^*=Q.
\]
All rearrangements are justified because the block index set is finite.
Thus \(N\) is self-adjoint, and
\eqref{shared-period-signed-orbital-bound} says
\(S^*S\le\alpha I_{\mathcal E}\).
Since \(S^*S\ge0\), its norm is at most \(\alpha\);
\(\|SS^*\|=\|S^*S\|\) gives \(Q\le\alpha I_H\).
For the truncation assertion, pass to the limit in the quadratic form
of each fixed \(v\). Taking self-adjoint parts of the truncations
does not change their real quadratic forms.
\end{proof}

The uniform bound \eqref{shared-period-signed-orbital-bound}, calibrated
so that \(\alpha D/M-1\) is sufficiently small, is an additional
arithmetic input. Bounds for the separate squared norms
\(\|K_{ij}\|_{\mathrm{HS}}^2\) do not by themselves establish this
signed operator estimate at that scale. Off-diagonal identity cells
are retained unless their vanishing has been proved separately.


Whenever the final matrices $A_i$ in this comparison represent the
actual arithmetic maps, Theorem~\ref{rh-centered-slope} may be applied
to them with full target dimension equal to $n$. Its target-cut
parameter is then zero; the retained mass and the final diagonal
and target energies carry the cost of projection. Establishing the
upper Bessel estimate and the same-map identification remains part
of the arithmetic construction.
